% =====================================================================
%  Chapter 9: Symanzik improvement and RG actions
%  Gattringer & Lang, "Quantum Chromodynamics on the Lattice" (LNP 788)
%  Reconstructed from extract/ch09.txt; equations auto-numbered.
% =====================================================================

\chapter{Symanzik improvement and RG actions}

When introducing the QCD action on the lattice we had to discretize the
derivative terms that show up in the continuum action. It was pointed out
then that any discretization, e.g., symmetric differences for the first
derivative in the fermion action, gives rise to discretization effects.
Typically the discretization effects are of $\order{a}$ for fermions and of
$\order{a^2}$ for the gauge fields. These disappear only in the continuum
limit when the lattice spacing $a$ is sent to zero. Performing the continuum
limit is, however, a nontrivial task. As one decreases $a$, the number of
lattice points has to increase, such that the physical volume remains constant
(ideally one would first send the number of lattice points to infinity before
sending $a$ to zero). Thus in a numerical simulation one always works with
finite $a$ and the discretization errors have to be dealt with, e.g., by
including them in the extrapolation to vanishing $a$.

An elegant way of approaching the problem is a systematical reduction of the
discretization errors. We have already mentioned that the discretization we
have chosen is not unique. Also other discretizations give rise to the same
formal continuum limit. In particular one may combine different terms to
obtain a lattice action with reduced discretization effects. Adding, e.g., an
extra term to the Wilson fermion action and matching its coefficient
appropriately, one can reduce the discretization error from $\order{a}$ to
$\order{a^2}$. In a similar way it is possible (and necessary for a full
improvement) to reduce also the discretization errors of the observables used.
A systematic implementation of these ideas is the Symanzik improvement program
which we discuss in Sect.~9.1.

A conceptionally different approach uses actions constructed from
renormalization group (RG) transformations. Improvement is achieved by
integrating out (blocking) short distance degrees of freedom and including
their effect in the action for the blocked fields. This strategy is presented
in Sects.~9.2--9.4. After discussing the blocking approach for free fermions
where the lattice action can be computed in closed form (Sect.~9.2), we
introduce the general setting for full QCD in Sect.~9.3. In Sect.~9.4 we
finally address an important use of the RG approach, the identification of the
lattice counterpart of continuum symmetries, and as an application we derive
the Ginsparg--Wilson equation.

\section{The Symanzik improvement program}
\label{sec:9.1}

\subsection{A toy example}
\label{sec:9.1.1}

Let us begin our discussion of improvement with a toy example which already
contains most of the steps that will be taken when improving lattice QCD. We
consider the (symmetric) discretization of the derivative $f'(x)$ for some
function $f(x)$ of a single real variable $x$:
\begin{equation}
  \frac{f(x+a) - f(x-a)}{2a}
  = f'(x) + a^2\,C^{(2)}(x) + a^4\,C^{(4)}(x) + O(a^6) .
\end{equation}
Note that due to the (anti-)symmetrical discretization on the left-hand side,
only even powers of $a$ and odd derivatives of $f$ can appear on the
right-hand side. We have introduced the abbreviations $C^{(2)}$ and $C^{(4)}$
for the correction terms. Since $a$ has dimension of length, in order to have
the same dimensions on both sides the correction terms $C^{(k)}$ must have
dimension $1+k$, i.e., $C^{(k)} \sim \mathrm{length}^{-(1+k)}$ (assuming that
$f$ is dimensionless).

Using the Taylor series expansion
\begin{equation}
  f(x \pm a) = f(x) \pm a\,f'(x) + \frac{a^2}{2}\,f''(x)
    \pm \frac{a^3}{6}\,f'''(x) + O(a^4) ,
\end{equation}
we can identify the leading correction term as
\begin{equation}
  C^{(2)}(x) = \frac{1}{6}\,f'''(x) .
\end{equation}
We stress at this point that this expression for the correction term is given
in continuum language, i.e., as a higher derivative of $f$.

The strategy for improvement is to add to the left-hand side of (9.1) a
discretized expression (i.e., an expression built from $f(x)$, $f(x\pm a)$,
$f(x\pm 2a)$, $\ldots$) such that the correction terms on the right-hand side
are canceled up to the requested order. For improvement of $\order{a^2}$ we
thus use the ansatz
\begin{equation}
  \frac{f(x+a) - f(x-a)}{2a} + c\,a^2\,D^{(3)}[f](x) = f'(x) + O(a^4) ,
\end{equation}
where $D^{(3)}[f]$ is a discretized expression obeying $D^{(3)}[f] \approx
f''' + O(a^2)$ and $c$ is some constant.

Employing again the power series (9.2), it is easy to see that
\begin{equation}
  D^{(3)}[f](x) = \frac{f(x+2a) - 2f(x+a) + 2f(x-a) - f(x-2a)}{2a^3} ,
  \qquad c = -\frac{1}{6}
\end{equation}
does the job and $\order{a^2}$ improvement is achieved. We remark, however,
that the choice (9.5) is not unique, and, e.g., also terms including
$f(x\pm 3a)$ could have been used.

Let us summarize the steps taken in our toy example, which already outline the
approach for improving lattice QCD:
\begin{itemize}
  \item We start from a simple discretized expression for the quantity of
        interest (the first derivative $f'$ in our example).
  \item Correction terms are identified using continuum language (higher
        derivatives in the toy example).
  \item The correction terms have certain symmetries (only odd derivatives in
        the example) and are ordered according to their dimension.
  \item In order to achieve improvement, discretized versions of the
        correction terms are added with suitable coefficients, such that
        corrections up to the desired order vanish.
  \item There is some arbitrariness in the choice of the discretized
        correction terms.
\end{itemize}

Exactly the same steps and features will appear in the improvement of lattice
QCD. The main difference is the determination of the coefficients. In the toy
example the coefficient $c$ followed from simple algebraic considerations. Due
to the nonlinear nature of QCD and the necessary renormalization, the
determination of the corresponding coefficients in QCD is much more involved
and must be done perturbatively or through a nonperturbative matching
procedure. The approach to improvement outlined here is known as Symanzik
improvement program [1--4].

\subsection{The framework for improving lattice QCD}
\label{sec:9.1.2}

When improving lattice QCD, we need to improve Euclidean correlation
functions, e.g., two point functions of interpolators $O_1$ and $O_2$ located
at two different space--time points $n, m \in \Lambda$,
\begin{equation}
  \ev{O_1(n)\,O_2(m)}
  = \frac{1}{Z}\int D[U]\,D[\psi,\bar\psi]\,
    e^{-S[U,\psi,\bar\psi]}\,O_1[U,\psi,\bar\psi;n]\,O_2[U,\psi,\bar\psi;m] ,
\end{equation}
which may be used to compute the energy spectrum and some matrix elements. To
obtain improvement for all these observables in general we will have to
improve both the lattice action $S$ and the interpolators $O_1$, $O_2$.

Let us begin the discussion with the improvement of the action. To be
specific, we start on the lattice with Wilson's gauge action (2.49) and a
fermion action with Wilson's Dirac operator (5.51). From the discussion of the
previous chapters we expect discretization errors of $\order{a}$ for the
fermion part and $\order{a^2}$ for the gauge part of the action.

Following the strategy illustrated in the toy example, we begin with
identifying a continuum expression for the correction terms. These should be
ordered according to their dimension and have the symmetries of the QCD
action. In other words we write down an effective action which describes the
behavior of Wilson's form of lattice QCD at finite $a$ and its approach to the
continuum limit. Following [1, 2, 5--7] we write the effective action in the
form
\begin{equation}
  S_{\mathrm{eff}} = \int d^4x\,
    \left[\mathcal{L}^{(0)}(x) + a\,\mathcal{L}^{(1)}(x)
    + a^2\,\mathcal{L}^{(2)}(x) + \ldots\right] .
\end{equation}
Here $\mathcal{L}^{(0)}$ is the usual QCD Lagrangian as defined in (2.3) and
(2.17). The terms $\mathcal{L}^{(k)}$, $k \geq 1$ are the additional
correction terms, which are built from products of quark and gluon fields such
that they have dimensions $4+k$, i.e., $\mathcal{L}^{(k)} \sim
\mathrm{length}^{-(4+k)}$. Thus, compared to $\mathcal{L}^{(0)}$, which is of
dimension 4, these terms will contain additional derivatives or powers of the
quark mass $m$. Requiring the symmetries of the lattice action (compare
Sect.~5.4), one may show that the leading correction term
$\mathcal{L}^{(1)}(x)$ can be written as a linear combination of the following
dimension-5 operators ($\sigma_{\mu\nu} \equiv [\gamma_\mu,\gamma_\nu]/2i$):
\begin{align}
  \mathcal{L}^{(1)}_1(x)
  &= \bar\psi(x)\,\sigma_{\mu\nu}F_{\mu\nu}(x)\,\psi(x) , \nonumber\\
  \mathcal{L}^{(1)}_2(x)
  &= \bar\psi(x)\overrightarrow{D}_\mu(x)\overrightarrow{D}_\mu(x)\psi(x)
   + \bar\psi(x)\overleftarrow{D}_\mu(x)\overleftarrow{D}_\mu(x)\psi(x) , \nonumber\\
  \mathcal{L}^{(1)}_3(x)
  &= m\,\tr[F_{\mu\nu}(x)F_{\mu\nu}(x)] , \nonumber\\
  \mathcal{L}^{(1)}_4(x)
  &= m\left[\bar\psi(x)\gamma_\mu\overrightarrow{D}_\mu(x)\psi(x)
   - \bar\psi(x)\gamma_\mu\overleftarrow{D}_\mu(x)\psi(x)\right] , \nonumber\\
  \mathcal{L}^{(1)}_5(x)
  &= m^2\,\bar\psi(x)\psi(x) .
\end{align}
This list of operators may be further reduced by using the field equation
$(\gamma_\mu D_\mu + m)\psi = 0$, which gives rise to the two relations
\begin{gather}
  \mathcal{L}^{(1)}_1 - \mathcal{L}^{(1)}_2 + 2\mathcal{L}^{(1)}_5 = 0 ,
  \nonumber\\
  \mathcal{L}^{(1)}_4 + 2\mathcal{L}^{(1)}_5 = 0 .
\end{gather}
These relations may be used to eliminate the terms $\mathcal{L}^{(1)}_2$ and
$\mathcal{L}^{(1)}_4$ from the set of operators that span the leading
correction, and it can be shown [6, 7] that this is true beyond tree level
(where the field equations hold). Thus it is sufficient to work with only the
terms $\mathcal{L}^{(1)}_1$, $\mathcal{L}^{(1)}_3$, and $\mathcal{L}^{(1)}_5$.
Two of these, $\mathcal{L}^{(1)}_3$ and $\mathcal{L}^{(1)}_5$, are terms that
up to some factor are already present in the original action, such that they
can be accounted for by a redefinition of the bare parameters $m$ and $g$ (or
$\beta = 6/g^2$).

Thus for $\order{a}$ improvement of the Wilson lattice action it is sufficient
to add the Pauli term $\mathcal{L}^{(1)}_1$ such that we obtain for the
improved action
\begin{equation}
  S_I = S_{\mathrm{Wilson}} + c_{sw}\,a^5\,\frac{1}{2}\sum_{\mu<\nu}\sum_{n\in\Lambda}
    \bar\psi(n)\,\sigma_{\mu\nu}\widetilde{F}_{\mu\nu}(n)\,\psi(n) .
\end{equation}
The real coefficient $c_{sw}$ is often referred to as Sheikholeslami--Wohlert
coefficient due to the authors of [5], where the improved action (9.10) was
written down for the first time. As outlined in the toy example,
$\widetilde{F}_{\mu\nu}$ has to be a discretized version of the corresponding
term in $\mathcal{L}^{(1)}_1$, i.e., a lattice form of the field strength
tensor. A convenient but not unique choice is

\begin{figure}[htbp]
  \centering
  \includegraphics[width=0.7\textwidth]{fig91.png}
  \caption*{Fig. 9.1. \emph{Graphical representation of the sum
    $Q_{\mu\nu}(n)$ of plaquettes in the $\mu$--$\nu$ plane used for the
    discretization of the field strength in (9.11)}.}
  \label{fig:91}
\end{figure}

\begin{equation}
  \widetilde{F}_{\mu\nu}(n) = \frac{-i}{8a^2}\left(Q_{\mu\nu}(n) - Q_{\nu\mu}(n)\right) ,
\end{equation}
where $Q_{\mu\nu}(n)$ is the sum of plaquettes $U_{\mu,\nu}(n)$ (compare
(2.48)) in the $\mu$--$\nu$ plane as shown in Fig.~9.1,
\begin{equation}
  Q_{\mu\nu}(n) \equiv U_{\mu,\nu}(n) + U_{\nu,-\mu}(n) + U_{-\mu,-\nu}(n) + U_{-\nu,\mu}(n) ,
\end{equation}
which is a discretization of the continuum field strength tensor. Due to the
shape of the terms which is reminiscent of a clover leaf, the last term in
(9.10) is often referred to as clover term or clover improvement.

At this point we need to add a couple of comments: We saw that for $\order{a}$
improvement it is sufficient to add a single term to the fermion action. The
only pure glue term $\mathcal{L}^{(1)}_3$ was found to be proportional to the
original gauge action and thus was absorbed in a redefinition of the bare
gauge coupling. Relevant purely gluonic operators appear only at dimension 6,
i.e., they contribute at $\order{a^2}$. This is in agreement with previous
remarks that for bosonic fields the discretization errors are only of
$\order{a^2}$. For the improvement of the gauge action we refer the reader to
the original literature, e.g., [3, 8], where the so-called L\"uscher--Weisz
gauge action is presented.

The second comment concerns the discretization errors of chirally symmetric
Dirac operators that obey the Ginsparg--Wilson equation (7.29). If one would
add the order $\order{a}$-improving term, i.e., the clover term, the resulting
Dirac operator would no longer obey the Ginsparg--Wilson equation. Thus we
must conclude that a Ginsparg--Wilson Dirac operator is already $\order{a}$
improved. As a matter of fact the nonlinear right-hand side of the
Ginsparg--Wilson equation (7.29) generates a lattice discretization of the
Pauli term $\mathcal{L}^{(1)}_1$ when the naive lattice Dirac operator (which
necessarily is a building block of any lattice Dirac operator) is squared.
This link between $\order{a}$ improvement and chiral symmetry hints already at
a nonperturbative strategy [6] for determining the Sheikholeslami--Wohlert
coefficient, which we will present below, following the discussion of the
improvement of currents.

\subsection{Improvement of interpolators}
\label{sec:9.1.3}

The improvement of the action, which we have discussed in the last section, is
already sufficient for $\order{a}$ improvement of on-shell quantities such as
hadron masses. However, for a full $\order{a}$ improvement of correlator
(9.6), which is necessary for the improvement of the hadronic matrix elements
computed from it, also the interpolators $O_i$ need to be improved.

As for the action the strategy is to find continuum expressions for the
correction terms which are organized with respect to their dimension and are
identified by requiring the symmetries of the operator $O_i$ one wants to
improve. Since currents with arbitrary quantum numbers may be constructed, it
is not possible to give an exhaustive list of terms needed for general
$\order{a}$ improvement of $n$-point functions. Following [6, 7] we instead
discuss two examples, the improvement of the isovector axial current $A_\mu^a$
and the pseudoscalar density $P^a$:
\begin{align}
  A^a_\mu(n) &= \frac{1}{2}\,\bar\psi(n)\gamma_\mu\gamma_5\tau^a\psi(n) , \nonumber\\
  P^a(n) &= \frac{1}{2}\,\bar\psi(n)\gamma_5\tau^a\psi(n) .
\end{align}
Here $\tau^a$ denotes one of the Pauli matrices $\tau^1$, $\tau^2$, $\tau^3$
acting in $\Nf = 2$-flavor space. The dimension-4 operators needed for
$\order{a}$ improvement of the current $A^a_\mu$ are given by
\begin{align}
  (A_1)^a_\mu(x) &= \frac{1}{2}\,\bar\psi(x)\gamma_5\sigma_{\mu\nu}
    \left[\overrightarrow{D}_\nu(x) - \overleftarrow{D}_\nu(x)\right]
    \tau^a\psi(x) , \nonumber\\
  (A_2)^a_\mu(x) &= \frac{1}{2}\,\partial_\mu\left[\bar\psi(x)\gamma_5\tau^a\psi(x)\right] , \nonumber\\
  (A_3)^a_\mu(x) &= \frac{m}{2}\,\bar\psi(x)\gamma_\mu\gamma_5\tau^a\psi(x) .
\end{align}
As for the improvement of the action we can apply the field equation, which
can be used to show that $A_1$ can be written as a linear combination of $A_2$
and $A_3$. Furthermore $A_3$ has the form of the original current of (9.13)
such that it can be absorbed in the overall multiplicative renormalization
constant $Z_A$ for the axial vector current (compare Chap.~11). Thus we obtain
for the $\order{a}$-improved axial current $A_I$ the expression
\begin{equation}
  (A_I)^a_\mu(n) = A^a_\mu(n) + c_A\,a\,\widetilde{\partial}_\mu\,P^a(n) ,
\end{equation}
where the real coefficient $c_A$ still has to be determined. By
$\widetilde{\partial}_\mu$ we denote the symmetric difference operator
$\widetilde{\partial}_\mu f(n) = (f(n + \hat\mu) - f(n - \hat\mu))/2a$.

For the pseudoscalar density $P^a$ the only term that occurs is a
normalization with an extra factor of the mass, which is again absorbed in the
renormalization. Thus for $\order{a}$ improvement it is sufficient to work
with $(P_I)^a = P^a$.

\subsection{Determination of improvement coefficients}
\label{sec:9.1.4}

For the $\order{a}$ improvement of $n$-point functions of the axial currents
we still need to determine the coefficients $c_{sw}$ and $c_A$. As already
remarked, here QCD deviates from our toy example due to its nonlinearity and
the necessary renormalization.

Using a traditional approach, one can compute the coefficients perturbatively
[5, 10, 11]:
\begin{equation}
  c_{sw} = 1 + 0.2659\,g^2 + O(g^4) , \qquad c_A = -0.00756\,g^2 + O(g^4) .
\end{equation}
However, for many important applications of lattice QCD one works in a regime
where the coupling $g$ is not small and the perturbative result is not
particularly useful. For an introductory discussion of another approach,
tadpole improvement [12], see [13].

We have already remarked that for the nonperturbative determination of $c_{sw}$
and $c_A$ one can explore the chiral symmetry of QCD [6]. In Chap.~11 we will
show that nontrivial relations between vacuum expectation values may be
derived, the so-called Ward identities. Here we simply state their existence
without proof. The particular Ward identity which we make use of here reads in
the continuum
\begin{equation}
  \ev{\partial_\mu A^a_\mu(x)\,O} = 2m\,\ev{P^a(x)\,O} .
\end{equation}
In this equation $A^a_\mu$ and $P^a$ are the continuum counterparts of the
corresponding lattice interpolators defined in (9.13). $O$ is an interpolator
chosen such that $\ev{P^a(x)\,O}$ does not vanish, with the restriction that
its support does not include the point $x$ where $A^a_\mu(x)$ and $P^a(x)$ are
located. Equation (9.17) expresses the fact that the isovector axial current
is conserved for a chiral theory where $m = 0$. Thus it is referred to as the
partially conserved axial current relation (PCAC).

On the lattice we expect that after renormalization (see Chap.~11), the
lattice interpolators approach the continuum relation (9.17), such that
\begin{equation}
  \ev{\widetilde{\partial}_\mu A^{(r)a}_\mu(n)\,O}
  = 2m^{(r)}\,\ev{P^{(r)a}(n)\,O} + O(a^z) .
\end{equation}
The superscript $(r)$ in this equation expresses the fact that renormalized
quantities are being used. The important observation is that the correction
$O(a^z)$ is different for the improved theory with properly chosen $c_{sw}$ and
$c_A$, where one expects $z = 2$, while for the unimproved theory one expects
$\order{a}$ corrections, i.e., $z = 1$. Thus we may use the corrections in the
PCAC relation (9.18) to determine $c_{sw}$ and $c_A$.

A slight complication arises through the necessary renormalization. Using the
improved axial current $(A_I)^a_\mu$ (for the pseudoscalar we work with the
original definition, $(P_I)^a = P^a$), we write the renormalized interpolators
as
\begin{align}
  A^{(r)a}_\mu(n) &= Z_A(1 + b_A\,a\,m)\left[A^a_\mu(n) + c_A\,a\,\widetilde{\partial}_\mu P^a(n)\right] , \nonumber\\
  P^{(r)a}(n) &= Z_P(1 + b_P\,a\,m)\,P^a(n) .
\end{align}
Here $Z_A$ and $Z_P$ are the renormalization constants for the axial vector
and the pseudoscalar, respectively (see Chap.~11). The real coefficients $b_A$
and $b_P$ parameterize their lattice corrections including correction terms
such as $(A_3)^a_\mu$, which, as we have already remarked, were absorbed in the
renormalization.

Based on (9.18) we define an unrenormalized quark mass $m_{\mathrm{AWI}}$,
referred to as PCAC mass or AWI (axial Ward identity) mass through the ratio
\begin{equation}
  m_{\mathrm{AWI}}(n)
  = \frac{\left\langle \left[\widetilde{\partial}_\mu A^a_\mu(n) + c_A\,a\,\Delta P^a(n)\right] O^a \right\rangle}
         {2\left\langle P^a(n)\,O^a \right\rangle} .
\end{equation}
In this equation $O^a$ is an interpolator which we will define below and
$\Delta$ denotes the Laplacian operator on the lattice, $\Delta f(n) =
\sum_\mu (f(n + \hat\mu) - 2f(n) + f(n - \hat\mu))/a^2$. From the lattice PCAC
relation (9.18) we obtain
\begin{equation}
  m_{\mathrm{AWI}} = \frac{Z_P(1 + b_P\,a\,m)}{Z_A(1 + b_A\,a\,m)}\,m^{(r)} + O(a^z) ,
\end{equation}
where $z = 2$ is obtained when $c_{sw}$ and $c_A$ have the correct
nonperturbative values. For the determination of these values one evaluates
$m_{\mathrm{AWI}}$ three times using three different $O^a$ in (9.20). This is
done at fixed bare parameters such that one of the three values can be used to
remove the ratio $(Z_P(1+b_P am))^{-1}(Z_A(1+b_A am))$ and the other two
determine $c_{sw}$ and $c_A$.

The remaining problem is the choice of the operators $O^a$ that one inserts
into (9.20). The operators $O^a$ must have the quantum numbers of $P^a$ and
should be chosen such that they probe the PCAC relation (and thus the cutoff
effect), but leave the renormalization factors in (9.21) unchanged. An
operator with that properties [6] may be constructed using the so-called
Schr\"odinger functional.

For the Schr\"odinger functional [14--18] one uses instead of periodic (or
antiperiodic) boundary conditions in time direction, Dirichlet, i.e., fixed
boundary conditions. For the spatial directions periodic boundary conditions
are kept. Using a lattice with $N_T + 1$ lattice points in time direction, we
need to specify the boundary values for the fields at $n_4 = 0$ and $n_4 =
N_T$:
\begin{align}
  P_+\psi(n,0) &= \rho(n) , & \bar\psi(n,0)P_- &= \bar\rho(n) , \nonumber\\
  P_-\psi(n,N_T) &= \rho'(n) , & \bar\psi(n,N_T)P_+ &= \bar\rho'(n) ,
\end{align}
where $P_\pm$ are the projectors $P_\pm = (1 \pm \gamma_4)/2$. For the gauge
fields we need to specify the values of the spatial gauge links at $n_4 = 0$
and $n_4 = N_T$:
\begin{align}
  U_j(n,0) &= \mathrm{diag}\left(e^{i\varphi_1/N}, e^{i\varphi_2/N}, e^{i\varphi_3/N}\right) , \nonumber\\
  U_j(n,N_T) &= \mathrm{diag}\left(e^{i\varphi'_1/N}, e^{i\varphi'_2/N}, e^{i\varphi'_3/N}\right) .
\end{align}
Here $N$ denotes the number of lattice points in spatial direction and the
phases $\varphi_i$, $\varphi'_i$ are subject to the conditions $\varphi_1 +
\varphi_2 + \varphi_3 = 0$ and $\varphi'_1 + \varphi'_2 + \varphi'_3 = 0$.

The dynamical variables, which are integrated over in the path integral
constituting the Schr\"odinger functional in (9.25) below, are just the
fermion and gauge degrees of freedom in the interior of our lattice. The
boundary values, on the other hand, may be used as source terms. Making use of
this feature we write down the interpolator $O^a$ which we insert in (9.21) as
\begin{equation}
  O^a = -\frac{a^6}{2}\sum_{n,m}\gamma_5\tau^a
    \frac{\partial}{\partial\rho(n)}\frac{\partial}{\partial\bar\rho(m)} .
\end{equation}
The vacuum expectation value of, e.g., the correlator of $A^a_\mu$ and $O^a$
in the Schr\"odinger functional is then defined as
\begin{equation}
  \ev{A^a_\mu(n)\,O^a}
  = \frac{1}{Z}\int D[\psi,\bar\psi,U]\,A^a_\mu(n)\,O^a\,e^{-S[\psi,\bar\psi,U]}
    \Big|_{\rho=\bar\rho=\rho'=\bar\rho'=0} ,
\end{equation}
i.e., the boundary values for the fermions are set to 0 after the derivatives
with respect to $\rho$, $\rho'$, $\bar\rho$, and $\bar\rho'$ in the definition
of $O^a$ have been performed. The boundary values for the gauge fields are
kept at nontrivial values and now may be used to probe the system.

The operator $O^a$ corresponds to a quark source and sinks with zero spatial
momentum placed at the timeslice $n_4 = 0$. Thus the vacuum expectation value
(9.25) may be interpreted as a quark and an antiquark lines traveling into the
interior of our lattice where they annihilate at the point $n$, i.e., the
point where we placed the current $A^a_\mu(n)$.

For our application the key property of the Schr\"odinger functional is that
in (9.21) the renormalization constants and the renormalized mass do not
depend on the boundary conditions $\varphi_i$ and $\varphi'_i$. Thus these can
be used to probe the exponent $z$ which parameterizes the discretization
error. If $c_{sw}$ and $c_A$ are chosen correctly the exponent is $z = 2$, and
the results for the PCAC mass $m_{\mathrm{AWI}}$ as defined in (9.20) should be
(essentially) independent of $\varphi_i$, $\varphi'_i$, and the coordinate
$n_4$ where the current $A^a_\mu$ is inserted. For technical details of the
strategy for determining $c_{sw}$ and $c_A$ from this condition we refer to the
original literature [19--21].

The numerical result for the coefficient $c_{sw}$ for two flavors of dynamical
Wilson fermions may be parameterized as [21]
\begin{equation}
  c_{sw} = \frac{1 - 0.454g^2 - 0.175g^4 + 0.012g^6 + 0.045g^8}{1 - 0.720g^2} ,
\end{equation}
which is a reliable approximation for the range of couplings $g \in [0.0,
1.1]$.

\section{Lattice actions for free fermions from RG transformations}
\label{sec:9.2}

In the last section we have achieved improvement by analyzing how different
terms that may appear in a lattice discretization of the action depend on the
lattice spacing $a$. We have seen that a linear combination of such terms with
suitably chosen coefficients reduces the discretization errors. In the next
three sections we now discuss a different approach, the so-called
renormalization group (RG) transformations, in which the degrees of freedom
above the lattice cutoff are integrated out and the effect of these UV modes
is incorporated in suitable terms of the lattice action.

In this section we directly block free fermions from the continuum (see [22]).
This warming-up exercise demonstrates some of the central concepts of RG
transformations for an example that can be solved in closed form. The example
is also of interest for the problem of chiral fermions on the lattice. We will
see that the lattice action we obtain from the blocking procedure obeys the
Ginsparg--Wilson equation, which we have already identified in Chap.~7 as the
key to chiral symmetry on the lattice.

\subsection{Integrating out the fields over hypercubes}
\label{sec:9.2.1}

The starting point is the continuum action for free massless fermions:
\begin{equation}
  S_F[\varphi,\bar\varphi] = \int d^4x\,\bar\varphi(x)\,\gamma_\mu\partial_\mu\,\varphi(x)
  = \int d^4p\,\bar\varphi(-p)\,i\gamma_\mu p_\mu\,\varphi(p) .
\end{equation}
For later use we have in the second step already transformed the action to
momentum space using the symmetric convention
\begin{equation}
  \bar f(p) = \int_{-\infty}^{\infty}\frac{d^4x}{(2\pi)^2}\,f(x)\,e^{-ip\cdot x} ,
  \qquad
  f(x) = \int_{-\infty}^{\infty}\frac{d^4p}{(2\pi)^2}\,\bar f(p)\,e^{ip\cdot x}
\end{equation}
for the Fourier transformation in the continuum. We stress at this point that
the Fourier transform in the continuum is denoted by $\bar f$, while for the
Fourier transform on the lattice $\hat f$ is used, as defined in Appendix~A.3.

The continuum fields $\varphi(x)$, $\bar\varphi(x)$ carry only Dirac indices,
since no gauge field is coupled. In order to obtain blocked lattice fields
$\varphi^B_n$, $\bar\varphi^B_n$ from the continuum fields $\varphi(x)$,
$\bar\varphi(x)$, we embed an infinite Euclidean lattice $\mathbb{Z}^4$ into
our Euclidean space--time $\mathbb{R}^4$. The lattice constant $a$ is set to
$a = 1$ for now, but will be displayed explicitly later when needed. We
construct the blocked fields by integrating the continuum fields over
hypercubes with sides of length 1, centered at the sites $n$ of our lattice,
i.e., $n \in \mathbb{Z}^4$:
\begin{equation}
  \varphi^B_n = \int_{-\infty}^{\infty}\frac{d^4p}{(2\pi)^2}
    \bar\varphi(p)\,\pi(p)\,e^{ip\cdot n} ,
  \qquad
  \pi(p) = \prod_{\mu=1}^4\int_{n_\mu-\frac{1}{2}}^{n_\mu+\frac{1}{2}}dx_\mu\,e^{ip\cdot x} .
\end{equation}
In this equation we have expressed the continuum field $\varphi(x)$ through
its Fourier transform $\bar\varphi(p)$ as defined in (9.28) and abbreviated
the remaining space--time integral as
\begin{equation}
  \pi(p) = \prod_{\mu=1}^4\int_{-\frac{1}{2}}^{\frac{1}{2}}dx_\mu\,e^{ip\cdot x}
  = \prod_{\mu=1}^4\frac{2\sin(p_\mu/2)}{p_\mu} .
\end{equation}
The function $\pi(p)$ describes the blocking procedure in momentum space. The
blocking of the conjugate field $\bar\varphi$ proceeds in exactly the same
way. Since the blocking procedure is trivial in Dirac space, the blocked
fields simply inherit the Dirac indices of their continuum counterparts.

Equation (9.29) expresses the blocked field $\varphi^B_n$ as a Fourier
integral with the Fourier transform $\bar\varphi(p)$ of the continuum. For
later use we rewrite (9.29) to obtain an alternative expression where the
lattice Fourier transform $\hat\varphi^B(p)$ of the blocked field is used:
\begin{equation}
  \varphi^B_n = \int_{-\pi}^{\pi}\frac{d^4p}{(2\pi)^4}\,\hat\varphi^B(p)\,e^{ip\cdot n} ,
  \qquad
  \hat\varphi^B(p) = \sum_{k\in\mathbb{Z}^4}\bar\varphi(p+2\pi k)\,\pi(p+2\pi k) .
\end{equation}
Note that now the momentum integral runs only over the Brillouin zone, i.e.,
$p_\mu \in (-\pi, \pi]$, which is the momentum space of an infinite lattice, as
can be seen from the formulas in Appendix~A.3, when one sends the number of
lattice points to infinity.

\subsection{The blocked lattice Dirac operator}
\label{sec:9.2.2}

After defining the blocked fields $\varphi^B_n$, $\bar\varphi^B_n$ and
evaluating their expression as Fourier integrals, we can now formally define
the lattice theory obtained from the blocking procedure:
\begin{multline}
  \exp\!\left[-\sum_{n,m}\bar\psi_n\,D_{nm}\,\psi_m\right]
  = \int D[\varphi,\bar\varphi]\,
    \exp\!\left[-\sum_n\left(\bar\psi_n-\bar\varphi^B_n\right)
      \left(\psi_n-\varphi^B_n\right)\right. \\
  \qquad\qquad \left.- S_F[\varphi,\bar\varphi]\right] .
\end{multline}
In this equation $\psi_m$, $\bar\psi_n$ denote the fields of the target
lattice theory, i.e., the fermions living on the infinite lattice
$\mathbb{Z}^4$. Consequently the indices $n$ and $m$ run over all of
$\mathbb{Z}^4$. The Dirac operator on the lattice is denoted by $D_{nm}$. The
left-hand side thus is the Boltzmann factor $\exp(-S_{\mathrm{latt}})$ of the
lattice theory. Both the lattice fields $\psi_m$, $\bar\psi_n$ and the lattice
Dirac operator $D_{nm}$ carry spinor indices which we suppress, i.e., we use
vector/matrix notation in Dirac space.

The right-hand side of (9.32) defines how the lattice theory, or more
precisely the lattice Dirac operator $D_{nm}$, is obtained from its continuum
counterpart: The continuum fields $\varphi$, $\bar\varphi$ are integrated over
in a path integral (we will comment on the continuum measure
$D[\varphi,\bar\varphi]$ later). In the Boltzmann factor we have the continuum
action $S_F[\varphi,\bar\varphi]$ from (9.27) as usual and an additional
quadratic piece which couples the blocked continuum fields $\varphi^B$,
$\bar\varphi^B$ to the fields $\psi$, $\bar\psi$ of the target theory on the
lattice.

Since the path integral on the right-hand side is Gaussian it can be solved in
closed form: It is convenient [22] to rewrite the blocking term with an
additional Gaussian integral over auxiliary Grassmann variables $\eta$,
$\bar\eta$:
\begin{multline}
  \exp\!\left[-\sum_n\left(\bar\psi_n-\bar\varphi^B_n\right)
    \left(\psi_n-\varphi^B_n\right)\right] \\
  = \int D[\eta,\bar\eta]\;
    \exp\!\left[\sum_n\left(\bar\eta_n\eta_n
      + \left(\bar\psi_n-\bar\varphi^B_n\right)\eta_n
      + \bar\eta_n\left(\psi_n-\varphi^B_n\right)\right)\right] .
\end{multline}
The auxiliary fields live on the lattice, carry a Dirac index, and the measure
$D[\eta,\bar\eta]$ is the product over the measures at all lattice points. The
terms in (9.33) that mix the auxiliary fields and the blocked fields
$\varphi^B$ are written as
\begin{align}
  \sum_n\bar\eta_n\,\varphi^B_n
  &= \int_{-\infty}^{\infty}\frac{d^4p}{(2\pi)^2}\,\bar J(-p)\,\bar\varphi(p) , \qquad
    \bar J(p) = \sum_n\bar\eta_n\,\pi(-p)\,e^{-ip\cdot n} , \nonumber\\
  \sum_n\bar\varphi^B_n\,\eta_n
  &= \int_{-\infty}^{\infty}\frac{d^4p}{(2\pi)^2}\,\bar\varphi(-p)\,J(p) , \qquad
    J(p) = \sum_n\eta_n\,\pi(-p)\,e^{-ip\cdot n} ,
\end{align}
where we have used (9.29) to express the blocked fields through their
continuum Fourier transforms.

Inserting (9.33) and (9.34) into the master formula (9.32), we obtain
\begin{equation}
  \exp\!\left[-\sum_{n,m}\bar\psi_n\,D_{nm}\,\psi_m\right]
  = \int D[\eta,\bar\eta]\,
    \exp\!\left[\sum_n\left(\bar\eta_n\eta_n + \bar\psi_n\eta_n + \bar\eta_n\psi_n\right)\right]
    I[\bar J,J] ,
\end{equation}
where we denote the remaining path integral over the continuum fields as
$I[\bar J,J]$ given by
\begin{equation}
  \begin{split}
    I[\bar J,J]
    &= \int D[\varphi,\bar\varphi]\,
      \exp\!\left[-\int_{-\infty}^{\infty}\frac{d^4p}{(2\pi)^2}\,
        \left(\bar\varphi(-p)\,i\gamma_\mu p_\mu\,\varphi(p)\right.\right. \\
    &\qquad\qquad\qquad
        \left.\left.+ \bar J(-p)\,\varphi(p) + \bar\varphi(p)\,J(p)\right)\right] \\
    &= C\,\exp\!\left[\int_{-\infty}^{\infty}\frac{d^4p}{(2\pi)^4}\,
        \bar J(-p)\,\frac{-i\gamma_\mu p_\mu}{p^2}\,J(p)\right] .
  \end{split}
\end{equation}
In the second step we have solved the path integral, as discussed in Chap.~5,
by completing the square in the exponent and a shift of the integration
variables in the path integral, $\varphi'(p) = \varphi(p) + i\gamma_\mu
p_\mu\,(2\pi)^{-2}p^{-2}J(p)$, and similarly for $\bar\varphi$. The remaining
term, which after the shift of the integration variables is independent of
$\bar J$ and $J$, is denoted by the constant $C$. We stress at this point that
we did not explicitly define the continuum measure $D[\varphi,\bar\varphi]$
and only assumed, that it remains invariant under a shift of variables. As a
matter of fact, for a proper definition of this measure, a regulator such as a
finite lattice is necessary. If one removes this regulator the constant $C$
diverges which, however, is irrelevant for the evaluation of vacuum
expectation values of operators, where such a constant cancels in the
normalization of expectation values with the partition function $Z$. Thus we
can ignore the constant $C$ and the path integral (9.36) is given by the term
quadratic in $J$ which comes from the completion of the square.

The final step is to insert result (9.36) into the remaining integral (9.35)
over $\eta$ and $\bar\eta$:
\begin{align}
  \exp\!\left[-\sum_{n,m}\bar\psi_n\,D_{nm}\,\psi_m\right]
  &= C\int D[\eta,\bar\eta]\;
     \exp\!\left[\sum_{n,m}\bar\eta_n\,(1+R)_{nm}\,\eta_m\right. \nonumber\\
  &\qquad \left.+ \sum_n\left(\bar\psi_n\eta_n + \bar\eta_n\psi_n\right)\right] \nonumber\\
  &= CC'\,\exp\!\left[-\sum_{n,m}\bar\psi_n\,(1+R)^{-1}_{nm}\,\psi_m\right] .
\end{align}
In the last step we have again solved the Gaussian path integral by completing
the square in the exponent and the remaining term, which is independent of
$\psi$ and $\bar\psi$, gives rise to another irrelevant constant $C'$. The
matrix $R_{nm}$ was obtained by inserting the expressions for $J(p)$ and
$\bar J(p)$ from (9.34) into the exponent on the right-hand side of (9.36).
Explicitly the matrix $R$ reads
\begin{equation}
  R_{nm} = \int_{-\infty}^{\infty}\frac{d^4p}{(2\pi)^2}\,
    \pi(-p)\,e^{-ip\cdot m}\,\frac{-i\gamma_\mu p_\mu}{p^2}\,\pi(p)\,e^{ip\cdot n} ,
\end{equation}
and by comparing the first and the last term in (9.37) we obtain the final
result for the lattice Dirac operator (after dropping the constants $C$ and
$C'$)
\begin{equation}
  D_{nm} = (1+R)^{-1}_{nm} .
\end{equation}
The lattice Dirac operator $D$ is most conveniently expressed by switching to
its Fourier transform:
\begin{equation}
  \hat D(p|q) = \frac{1}{(2\pi)^4}\sum_{n,m}e^{-ip\cdot n}\,D_{nm}\,e^{iq\cdot m} .
\end{equation}
For the necessary Fourier transform of $R$ one finds
\begin{align}
  \hat R(p|q)
  &= \delta(p-q)\sum_{k\in\mathbb{Z}^4}
     \frac{\pi(q+2\pi k)^2}{(q+2\pi k)^2}\,(-i)\gamma_\mu\,(q+2\pi k)_\mu \nonumber\\
  &= -i\,\delta(p-q)\,\gamma_\mu\,v_\mu(q) , \nonumber\\
  \text{with}\quad v_\mu(q)
  &= \sum_{k\in\mathbb{Z}^4}
     \frac{\pi(q+2\pi k)^2}{(q+2\pi k)^2}\,(q+2\pi k)_\mu .
\end{align}
Inserting this into (9.39) one obtains [22]
\begin{equation}
  \hat D(p|q) = \delta(p-q)\,\hat D(q) , \qquad
  \hat D(q) = \left(a\,\mathbb{1} - i\gamma_\mu v_\mu(q)\right)^{-1}
    = \frac{a\,\mathbb{1} + i\gamma_\mu v_\mu(q)}{a^2 + v(q)^2} .
\end{equation}
Here we have split off the Dirac-$\delta$ that ensures momentum conservation
which is a consequence of translation invariance. We have also re-inserted the
lattice spacing $a$ necessary for the analysis below.

\subsection{Properties of the blocked action}
\label{sec:9.2.3}

The lattice propagator in momentum space, $\hat D(q)$, is the result of our
blocking procedure for obtaining the effective lattice theory. It has a series
of remarkable properties which we discuss now.

First we convince ourselves that $\hat D(q)$ has the correct naive continuum
limit and that it is free of doublers. In other words we need to show that
\begin{equation}
  \hat D(q) = \begin{cases}
    i\gamma_\mu q_\mu + O(a) & \text{when all } q_\mu \approx 0 ,\\
    O(1/a) & \text{when at least one } q_\mu \approx \pi/a .
  \end{cases}
\end{equation}
When inspecting the final result (9.42) for the lattice Dirac operator, it is
obvious that the properties (9.43) must come from the behavior of the
functions $v_\mu(q)$ given in (9.41). As a matter of fact one can show (e.g.,
most quickly by analyzing the $v_\mu(q)$ numerically) that they behave as
\begin{equation}
  v_\mu(q) \approx \begin{cases}
    q_\mu/q^2 & \text{when all } q_\mu \approx 0 ,\\
    0 & \text{when at least one } q_\mu \approx \pi/a .
  \end{cases}
\end{equation}
This behavior exactly implies the required properties of a lattice Dirac
operator as stated in (9.43).

Having established that the blocked Dirac operator $\hat D$ gives rise to a
proper lattice theory let us now investigate its properties further. When
transforming back to real space it is found that the blocked Dirac operator is
not ultralocal, in agreement with the general discussion in Sect.~7.4.2. It
may, however, be demonstrated [22] that the blocked Dirac operator falls off
exponentially as is necessary for a local quantum field theory.

Another important aspect is discretization errors. We have seen in (9.43) that
in the limit of vanishing momentum $q$ we obtain the correct behavior
$i\gamma_\mu q_\mu$. What, however, are the deviations from that ideal behavior
as one increases the momentum? A physically interesting way of assessing this
question is to analyze the dispersion relation of a free fermion described by
$\hat D$. This may be done by computing the poles of $\hat D(q)^{-1}$ in the
complex $q_4$ plane for different values of the spatial momentum $q$. Plotting
the energy $E(q)$, which is given by the modulus of the position $q_4$ of the
pole, versus the spatial momentum, one finds that for the blocked Dirac
operator the values fall exactly on top of the linear continuum dispersion
relation, $E(q) = |q|$, essentially almost all the way up to the cutoff
$\pi/a$. This indicates that the discretization errors are small.

Finally we comment on the chiral properties of the blocked lattice Dirac
operator $\hat D(q)$. It is a simple exercise to show that it obeys the
Ginsparg--Wilson equation in the form
\begin{equation}
  \gamma_5\hat D(q) + \hat D(q)\gamma_5 = 2a\,\hat D(q)\gamma_5\hat D(q) .
\end{equation}
We remark that here we encounter the Ginsparg--Wilson equation with a slightly
different normalization than the original version (7.29), namely an extra
factor of 2 on the right-hand side. This factor is, however, completely
irrelevant, since the whole program for chiral symmetry on the lattice, which
we presented in Chap.~7, can be implemented with the modified normalization as
well. Thus we find that the blocked theory is also chiral.

To summarize this section, we have shown that blocking free fermions from the
continuum gives rise to
\begin{itemize}
  \item Proper lattice fermions without doublers and a local action.
  \item Fermions with a perfect dispersion relation almost up to the cutoff.
  \item A Dirac operator that obeys the Ginsparg--Wilson equation.
\end{itemize}
With these three properties it is obvious that the blocking procedure is a
powerful method for constructing lattice discretizations of quantum field
theories. In Sect.~9.4 we will show that the emerging chiral lattice symmetry
is not a coincidence but is the consequence of a rather general mapping of
continuum symmetries onto a corresponding lattice theory obtained by blocking.
Before we come to this analysis, we need, however, to address the problem of
generalizing the strategy, which we have here outlined for free fermions, to
the case of full QCD where gluons and quark fields are interacting.

\section{Real space renormalization group for QCD}
\label{sec:9.3}

In the last section we have demonstrated by an explicit calculation for free
fermions that the RG approach of integrating out degrees of freedom above the
cutoff $\sim 1/a$ is a powerful tool for constructing a lattice discretization
of a quantum field theory. The resulting lattice action has beautiful
properties, in particular discretization effects are absent almost up to the
cutoff and the lattice Dirac operator obeys the Ginsparg--Wilson equation,
which is a lattice manifestation of the continuum chiral symmetry.

For the case of full QCD, i.e., when gauge fields are coupled to the fermions,
the situation is considerably more complicated. A serious obstacle prevents a
direct blocking from the continuum like for free fermions considered in the
last section: The continuum path integral on the right-hand side of (9.32)
cannot be solved in closed form if non-Gaussian terms appear. Such terms are
the interaction terms of the gauge field and the fermions, as well as the
cubic and quartic terms of the gauge action.

The conceptual way out is to use a lattice version of QCD on a very fine
lattice, i.e., a lattice close to the continuum limit where the lattice
spacing $a$ is very small and discretization effects are tiny. This fine
lattice replaces the continuum theory as a starting point for the blocking
transformation. Instead of integrating hypercubes of the continuum we average
the lattice fields on a block of sites of the fine lattice to construct the
fields on a coarser lattice. The action on the coarse lattice is identified
through an equation equivalent to (9.32). Iterating such a real space
renormalization group transformation we may obtain the action of the target
theory on a lattice with the coarseness we want to work at. Up to the cutoff
this coarse action will have the same discretization errors as the theory on
the fine lattice we started with. The physics of the short-distance UV modes,
which are averaged over in the RG transformation, is encoded in terms of the
action used on the coarse lattice. We remark that originally the real space
renormalization group was developed as an approach for analyzing systems of
statistical mechanics near criticality [23, 24]. As an introduction to its
application in lattice QCD which goes beyond our brief sketch we recommend
[25].

\begin{figure}[htbp]
  \centering
  \includegraphics[width=0.85\textwidth]{fig92.png}
  \caption*{Fig. 9.2. \emph{Example of an overlapping discrete block spin
    transformation. We show for a plane of the lattice how the fermion fields
    on the fine lattice (filled dots) may be averaged to construct the blocked
    fermion fields living on the coarse lattice (open circles around the sites
    of the lattice with only even indices). Along the arrows gauge
    transporters are included to obtain gauge covariant spinors on the coarse
    lattice}.}
  \label{fig:92}
\end{figure}

\subsection{Blocking full QCD}
\label{sec:9.3.1}

Let us begin with describing how we can construct new fields living on a
coarse lattice $\Lambda'$ with spacing $a' = 2a$ from fields on a fine lattice
$\Lambda$ with lattice constant $a$. Since we now consider QCD, we have to
block both the fermions and the gauge fields. Before we discuss the actual
construction, we set the notational conventions: The space--time index of the
fine lattice $\Lambda$ is denoted by $n$ and the space--time index of the
coarse lattice $\Lambda'$ by $n'$. The fermion and the gauge fields on the
fine lattice are $\psi_n$, $\bar\psi_n$ and $U_\mu(n)$, their blocked
counterparts living on the coarse lattice are $\psi^B_{n'}$, $\bar\psi^B_{n'}$
and $U^B_\mu(n')$, and the genuine coarse fields are denoted by $\psi'_{n'}$,
$\bar\psi'_{n'}$ and $U'_\mu(n')$.

We start with discussing the blocking prescription of the fermions. Many
different blocking schemes are possible. A rather simple one is illustrated in
Fig.~9.2, where we show how the fields may be blocked in a plane of the
lattice (the other planes are blocked equivalently). The fermion fields
$\psi_n$, $\bar\psi_n$ live on the sites of the fine lattice (filled dots in
the figure) and have to be averaged suitably for obtaining the blocked
fermions $\psi^B_{n'}$, $\bar\psi^B_{n'}$, living on the coarse lattice which
we choose to consist of the sites with only even indices of the fine lattice
(marked by circles in the plot). The blocked fields $\psi^B_{n'}$,
$\bar\psi^B_{n'}$ at a site $n'$ of the coarse lattice are obtained by
averaging over all the fields $\psi_n$, $\bar\psi_n$ on all lattice sites $n$
which are connected to $n'$ by one or more oriented paths. Obviously the sites
of the fine lattice are connected to several sites of the coarse lattice. Such
a type of blocking prescription is referred to as overlapping block spin
transformation.

It is mandatory to keep gauge covariance intact when combining the fermion
fields $\psi_n$, $\bar\psi_n$ from different sites $n$ of the fine lattice.
Thus we need to connect the fields at sites $n$ to the lattice site $n'$ where
we put the blocked field through gauge transporters $U(n',n)$. Such
transporters are obtained by multiplying the link variables $U_\mu(m)$ of the
fine lattice along the oriented paths in the figure. If a link is run through
in negative direction we apply our usual convention $U_{-\mu}(m) \equiv
U_\mu(m-\hat\mu)^\dagger$. In order to obtain a blocking prescription which
respects the symmetries of the lattice, several sites $n$ are connected to
$n'$ by a set of paths which are related by symmetries such as discrete
rotations, reflections etc. We can summarize the blocking of the fermions as
\begin{equation}
  \psi^B_{n'} = \sum_{n\in\mathcal{N}(n')}
    \left(\sum_{p\in\mathcal{P}(n',n)}s_p\prod_{(m,\mu)\in p}U_\mu(m)\right)\psi_n
  \equiv \sum_n\omega[U]_{n'n}\,\psi_n .
\end{equation}
Here $\mathcal{N}(n')$ is a neighborhood of sites $n$ of the fine lattice
around the target site $n'$ on the coarse lattice. By $\mathcal{P}(n',n)$ we
denote the set of paths $p$ that we want to use for connecting the site $n$ of
the fine lattice to $n'$, e.g., the choice depicted in Fig.~9.2. The product
runs over all the links $(m,\mu)$ in the path $p$. Each path comes with a
coefficient $s_p$. These coefficients must reflect the symmetries (e.g., paths
related to each other by a rotation must have the same coefficient), and as a
normalization condition, we require the coefficients of all paths connected to
a target site $n'$ to sum up to 1. In the last step of (9.46) we have
rewritten the blocking prescription by introducing a matrix $\omega[U]_{n'n}$
which depends on the gauge fields $U$ of the fine lattice and collects all the
terms that connect $n'$ and $n$. The conjugate fields $\bar\psi(n)$ are
blocked equivalently with the hermitian conjugate of $\omega[U]$. We point out
that, as in the last section, the blocking procedure is trivial in Dirac space.

Concerning the blocking of the gauge link $U_\mu(n)$ one could simply multiply
the two fine link variables that connect two sites of the coarse lattice in
order to obtain the blocked variables. However, also here usually a more
general gauge covariant averaging is used, with a simple version given by
\begin{align}
  U^B_\mu(n') &= (1-\alpha)\,U_\mu(n')\,U_\mu(n'+\hat\mu) \nonumber\\
    &\quad + \frac{\alpha}{6}\sum_{\nu\neq\mu}
      \left[U_\nu(n')\,U_\mu(n'+\hat\nu)\,U_\mu(n'+\hat\nu+\hat\mu)\,
       U_\nu(n'+2\hat\mu)^\dagger \right. \nonumber\\
    &\qquad \left. + U_\nu(n'-\hat\nu)^\dagger\,U_\mu(n'-\hat\nu)\,
       U_\mu(n'-\hat\nu+\hat\mu)\,U_\nu(n'-\hat\nu+2\hat\mu)\right] .
\end{align}
Here $\hat\mu$ and $\hat\nu$ are unit vectors on the fine lattice pointing in
the $\mu$ and $\nu$ directions, and $\alpha$ is a real parameter of the
blocking procedure. Essentially the blocking linearly combines paths of links
on the fine lattice which connect neighboring sites on the coarse lattice.
Along these paths the ordered products of the fine link variables are taken
into account. Finding an optimal blocking prescription requires a certain
amount of experimenting, and prescriptions more refined than the one in (9.47)
are discussed in [26--29].

We remark that in the above definition the blocked field $U^B_\mu(n')$ is not
an element of the gauge group, since the sum of two SU(3) matrices is not in
SU(3). Below we will define how the blocked field $U^B_\mu(n')$ is connected
to the link variables $U'_\mu(n')$ on the coarse lattice, which of course will
then be elements of the gauge group.

Having defined how to obtain the blocked fields $\psi^B$, $\bar\psi^B$, and
$U^B$, we can now formulate a blocking relation which generalizes (9.32) and
defines the theory on the coarse lattice:
\begin{multline}
  e^{-S'_F[\psi',\bar\psi',U'] - \beta' S'_G[U']}
  = \int D[\psi,\bar\psi]\,D[U]\,e^{-S_F[\psi,\bar\psi,U] - \beta S_G[U]} \\
  \times e^{-T_F[\psi',\bar\psi',\psi^B,\bar\psi^B,U] - \beta\,T_G[U',U^B]} .
\end{multline}
In this equation $S_F[\psi,\bar\psi,U]$ and $S'_F[\psi',\bar\psi',U']$ are the
fermion actions for the fine and the coarse lattices, and $\beta S_G[U]$ and
$\beta' S'_G[U']$ the corresponding gauge actions, where $\beta$, $\beta'$
denote the inverse gauge couplings for the two lattices. The right-hand side
is now a well-defined lattice path integral over the degrees of freedom on the
fine lattice. The blocking functions $T_F$ and $T_G$ connect the fermion and
gauge fields $\psi'$, $\bar\psi'$, $U'$ of the target theory on the coarse
lattice to the blocked quantities $\psi^B$, $\bar\psi^B$, $U^B$. For the
fermions we use
\begin{align}
  T_F[\psi',\bar\psi',\psi^B,\bar\psi^B,U]
  &= \kappa_F\sum_{n'}\left(\bar\psi'_{n'} - b\,\bar\psi^B_{n'}\right)
     \left(\psi'_{n'} - b\,\psi^B_{n'}\right) \nonumber\\
  &= \kappa_F\sum_{n'}\left(\bar\psi'_{n'} - b\sum_n\bar\psi_n\,\omega[U]^\dagger_{nn'}\right) \nonumber\\
  &\qquad \times\left(\psi'_{n'} - b\sum_n\omega[U]_{n'n}\,\psi_n\right) ,
\end{align}
where $b$ is a real-valued normalization constant and $\kappa_F$ a real
parameter that may be used for tuning properties of the action on the coarse
lattice, in particular its locality. In the second line of (9.49) we have
expressed the blocked fields with the transformation matrix $\omega[U]$
introduced in (9.46). For the gauge field we define
\begin{align}
  T_G[U',U^B] &= \frac{\kappa_G}{3}\sum_{n',\nu}
    \left[\mathrm{Re}\,\tr\left(U'_\nu(n')\,U^B_\nu(n')^\dagger\right)
      - \mathcal{N}\left(U^B_\nu(n')\right)\right] , \nonumber\\
  \exp\left[\mathcal{N}\left(U^B_\nu(n')\right)\right]
  &\equiv \int dW\,\exp\left[\beta\kappa_G\,\mathrm{Re}\,\tr\left(W\,U^B_\nu(n')^\dagger\right)/3\right] .
\end{align}
The second equation defines the normalization $\mathcal{N}(U^B_\nu(n'))$
through an integral over $W \in \mathrm{SU}(3)$, and $\kappa_G$ is again a
real parameter of the transformation.

\subsection{The RG flow of the couplings}
\label{sec:9.3.2}

In the last section we have defined the blocking prescriptions for fermions
and gauge fields and thus have given a well-defined meaning to the
renormalization group equation (9.48) which connects the actions for the
theories on the coarse and the fine lattices. For an arbitrary coarse
configuration $\psi'$, $\bar\psi'$, $U'$, the right-hand side of (9.48)
defines the action on the coarse lattice in the exponent on the left-hand
side.

We have already remarked that the key idea of the RG transformation is to
include the physics of the modes on the fine lattice in suitable terms in the
action used for the coarse lattice. Such terms can, e.g., look like the clover
term (9.10) introduced in the first section of this chapter. All these terms
of the action on the coarse lattice come with some coefficients $c'_1$, $c'_2$
$\ldots$, which for the example of the clover term correspond to the
Sheikholeslami--Wohlert coefficient of the Symanzik improvement program. We
assume that we can view the RG transformation (9.48) as a mapping in the space
of the couplings which parameterize a lattice QCD action,
\begin{equation}
  \beta,\,c_1,\,c_2,\,\ldots\;\xrightarrow{\;\mathrm{RG}\;}
    \beta',\,c'_1,\,c'_2,\,\ldots ,
\end{equation}
where the $\beta$, $c_i$ are the couplings on the fine lattice (most of which
vanish for the plain Wilson action which might be the action we start from).
For labeling the couplings we do not distinguish between the fermion and gauge
field parts of the action and only single out the inverse gauge couplings
$\beta$, $\beta'$.

Starting from the theory with couplings $\beta$, $c_1$, $c_2$ $\ldots$, we
block the fields and the couplings of the coarse lattice, $\beta'$, $c'_1$,
$c'_2$, $\ldots$, assume values such that the effects of the integrated-out
fields of the fine lattice are taken into account. At the same time the
lattice spacing increases by a scale factor of 2,
\begin{equation}
  a \xrightarrow{\;\mathrm{RG}\;} a' = 2a .
\end{equation}
We remark that the scale factor of 2 is particular for the blocking
transformation defined by Fig.~9.2 and (9.47), and other blocking
prescriptions can have a different scale factor.

We now ask ourselves whether mapping (9.51) has a fixed point (FP), i.e., a
set of couplings that is mapped onto itself under the RG transformation
(9.51). Equation (9.52) suggests that two types of fixed points might exist,
corresponding to $a = a' = \infty$ or to $a = a' = 0$. Obviously only the
second possibility, which often is referred to as critical fixed point, is of
interest to us since it corresponds to the continuum limit $a \to 0$. As we
have discussed in Sect.~3.5.4, the continuum limit of lattice QCD, where $a
\to 0$, is obtained in the limit $\beta \to \infty$. Thus we expect to find a
suitable critical fixed point at $\beta = \infty$ and parameter values $c_1$,
$c_2$, $\ldots$ . Under our RG transformation (9.51) these values are mapped
onto themselves.

In a practical application we cannot start our RG transformation from $\beta =
\infty$. However, a detailed analysis of the flow of couplings under repeated
RG transformations (see, e.g., [25]) suggests that it is sufficient to start
from a large but finite value of $\beta$ and a few steps of repeated RG
transformations will drive the couplings $c_1$, $c_2$, $\ldots$ into the fixed
point values $c_1$, $c_2$, $\ldots$ .

\subsection{Saddle point analysis of the RG equation}
\label{sec:9.3.3}

So far we have specified the blocking prescription for full QCD, defined the
RG transformation (9.48), and identified a possible fixed point which
corresponds to the continuum physics we are interested in. What is left to do
is to find a way of solving the RG transformation (9.48), such that we can
construct the RG action on the lattice.

For this purpose we may use the fact that the critical fixed point we want to
find corresponds to $\beta = \beta' = \infty$. In this limit the path integral
on the right-hand side of (9.48) is dominated by its saddle point, i.e., a
configuration (not necessarily unique) of gauge fields $U$ on the fine lattice
which minimizes the exponent. On both sides of (9.48) the terms in the
exponent which are multiplied with the inverse couplings $\beta$, $\beta'$
have to match as we send the inverse couplings to infinity. We read off the
following equation for the gauge field part of the action:
\begin{equation}
  S'_G[U'] = \min_U\left\{S_G[U] + T^\infty_G[U',U^B]\right\} .
\end{equation}
The superscript $\infty$ attached to the blocking function $T_G$ indicates
that it also has to be evaluated at $\beta = \infty$. Since we want to identify
the so-called fixed point (FP) action $S^{\mathrm{FP}}_G$ at the fixed point,
which is defined by the coefficients $c_i$ that are mapped onto themselves, we
can set $S'_G = S_G = S^{\mathrm{FP}}_G$ and obtain the equation that
identifies the gauge part $S^{\mathrm{FP}}_G$ of the fixed point action,
\begin{align}
  S^{\mathrm{FP}}_G[U']
  &= \min_U\left\{S^{\mathrm{FP}}_G[U] + T^\infty_G[U',U^B]\right\} \nonumber\\
  &= \min_U\Biggl\{S^{\mathrm{FP}}_G[U] + \frac{\kappa_G}{3}\sum_{n',\nu}
     \left[\mathrm{Re}\,\tr\left(U'_\nu(n')\,U^B_\nu(n')^\dagger\right)\right. \nonumber\\
  &\qquad \left.- \mathcal{N}^\infty\left(U^B_\nu(n')\right)\right]\Biggr\} ,
\end{align}
where in the second step we have inserted the explicit expression of the
blocking function $T_G$ as given in (9.50). The necessary limit $\beta \to
\infty$ is easy to perform since $\beta$ enters only through the normalization
constant $\mathcal{N}^\infty$ which again can be evaluated through its saddle
point giving
\begin{equation}
  \mathcal{N}^\infty\left(U^B_\nu(n')\right)
  = \frac{\kappa_G}{3}\,\max_{W\in\mathrm{SU}(3)}
    \mathrm{Re}\,\tr\left(W\,U^B_\nu(n')^\dagger\right) .
\end{equation}
Matching the exponents with factors $\beta$, $\beta'$ in the saddle point
analysis of (9.48) leads to Eq.~(9.54) which determines the gauge part of the
fixed point action. The requirement that also the remaining factors from the
fermion contributions match gives rise to
\begin{equation}
  e^{-S'_F[\psi',\bar\psi',U']} = \int D[\psi,\bar\psi]\,
    e^{-S_F[\psi,\bar\psi,U^{\min}] - T_F[\psi',\bar\psi',\psi^B,\bar\psi^B,U^{\min}]} ,
\end{equation}
where the right-hand side has to be evaluated for the configuration $U^{\min}$
which minimizes the right-hand side of the gauge field RG equation (9.54).

The integral on the right-hand side of (9.56) can be solved, since it is
Gaussian. We write the coarse and fine fermion actions with the help of Dirac
operators $D'$ and $D$ for the coarse and the fine lattices as
\begin{align}
  S'_F[\psi',\bar\psi',U'] &= \sum_{n',m'}\bar\psi'_{n'}\,D'[U']_{n'm'}\,\psi'_{m'} , \nonumber\\
  S_F[\psi,\bar\psi,U^{\min}] &= \sum_{n,m}\bar\psi_n\,D[U^{\min}]_{nm}\,\psi_m .
\end{align}
Using the fact that the blocking function $T_F(\psi',\bar\psi',\psi^B,\bar\psi^B,
U^{\min})$ is also bilinear in $\psi$ and $\bar\psi$ (see (9.49)) we solve the
fermionic integral with the techniques of Sect.~5.1.4. The result of the
integral on the right-hand side of (9.56) is the exponential of a bilinear
form for the coarse fields $\psi'$, $\bar\psi'$. The kernel of this bilinear
form has to match the kernel in the exponent of the bilinear form on the
left-hand side of (9.56). We obtain
\begin{equation}
  \begin{split}
    D'[U']_{n'm'}
    &= \kappa_F\,\delta_{n'm'}
       - \kappa_F^2 b^2\left[\omega[U^{\min}]
         \left(D[U^{\min}] + \kappa_F b^2\,\omega[U^{\min}]^\dagger
         \omega[U^{\min}]\right)^{-1}\right. \\
    &\qquad\qquad \left.\cdot \omega[U^{\min}]^\dagger\right]_{n'm'} ,
  \end{split}
\end{equation}
which is the RG equation that relates the Dirac operator $D'$ on the coarse
lattice to the Dirac operator $D$ of the fine lattice.

Equations (9.54) and (9.58) are the two RG equations that connect the action
on the fine lattice to the action on the coarse lattice. It is interesting to
note that the RG equation for the fixed point gauge action is independent of
the problem for the fermions. The RG equation (9.58) for the fermions,
however, needs as input a gauge configuration $U^{\min}$ on the fine lattice,
which minimizes the right-hand side of (9.54). In the next section we discuss
how solutions of (9.54) and (9.58) may be approximated.

\subsection{Solving the RG equations}
\label{sec:9.3.4}

The two RG equations (9.54) and (9.58) are nonlinear equations which should
hold for arbitrary coarse background gauge fields $U'$. In a first step one has
to determine the fine configuration $U^{\min}$ that corresponds to a given
$U'$ and minimizes the right-hand side of (9.54). This configuration is then
used in the RG equation (9.58) for the fermions.

In order to determine the fixed point gauge action $S^{\mathrm{FP}}_G$ and the
Dirac operator $D'$, both these quantities have to be approximated and
parameterized using a finite number of terms with couplings $c_i$. For the
gauge action one may use a linear combination of traces over closed loops of
link variables and powers thereof,
\begin{equation}
  S^{\mathrm{FP}}_G[U] = \sum_m\sum_{l\in\mathcal{L}}c_m(l)\,
    \left(\frac{1}{3}\,\mathrm{Re}\,\tr[1-U_l]\right)^m .
\end{equation}
In this equation $\mathcal{L}$ is a collection of closed loops on the lattice
and $U_l$ the ordered product of link variables along such a loop $l \in
\mathcal{L}$. The couplings $c_m(l)$ are real parameters of the action. The set
$\mathcal{L}$ of loops has to be chosen such that the lattice symmetries,
e.g., discrete translations, rotations, are obeyed. Obviously, the
parameterized action (9.59) is a generalization of the Wilson gauge action
(2.49) and reduces to this action if we choose for $\mathcal{L}$ the set of all
plaquettes and restrict ourselves only to powers $m = 1$ in (9.59).

The parameterization of the Dirac operator proceeds in a similar way [30, 31]
and may also be viewed as a generalization of the Wilson Dirac operator (5.51).
When we discussed the doubling problem in Sect.~5.2 we were forced to include
terms that come with the unit matrix in Dirac space. In the ansatz for the
parameterized Dirac operator now all possible 16 generators $\Gamma_k \in \{1,
\gamma_\mu, \sigma_{\mu\nu}, \gamma_\mu\gamma_5, \gamma_5\}$ of the Clifford
algebra are used. Similar to the gauge action we allow for a whole collection
of paths $p \in \mathcal{P}^{(k)}_{n,m}$ that connect two sites $n$, $m$ on the
lattice and we obtain [30, 31]
\begin{equation}
  D[U]_{nm} = \sum_{k=1}^{16}\Gamma_k
    \sum_{p\in\mathcal{P}^{(k)}_{n,m}}c^{(k,p)}_{n,m}[U]\,U_p ,
\end{equation}
where $U_p$ denotes the product of link variables along the path $p$. Again
the set of paths $\mathcal{P}^{(k)}_{n,m}$ has to respect the discrete lattice
symmetries. In addition the set of paths $\mathcal{P}^{(k)}_{n,m}$ and the
coefficients $c^{(k,p)}_{n,m}$ depend on the element of the Clifford algebra,
since for the fermions the reflections and other symmetries such as charge
conjugation (see Sect.~5.4) involve also matrices acting in Dirac space. The
coefficients $c^{(k,p)}_{n,m}[U]$ may be gauge-invariant functions of the link
variables $U_\mu$, e.g., they could depend on the local plaquette values.

For the gauge action $S^{\mathrm{FP}}_G$ the couplings on both sides of (9.54)
are the same and determine the fixed point action which then is used at
arbitrary $\beta$. The actual determination is done via minimizing the
difference between the left-hand and the right-hand sides of (9.54). This step
actually involves two minimizations: The coefficients of the action and the
fine configuration have to be varied. In a practical approach [26--29, 32] one
first uses a simpler gauge action on the right-hand side of (9.54) which
serves to determine a minimal gauge configuration $U^{\min}$ corresponding to
a coarse configuration $U'$. In a second step this minimal configuration is
then used to determine the coefficients $c_m(l)$ of the parameterized fixed
point action (9.59). The resulting fixed point gauge action was tested
nonperturbatively and was shown to have only very small cutoff effects
[26--29].

For the fermionic RG equation (9.58) an approach, slightly different from the
fixed point action for the gauge fields, was chosen [33--37]. Instead of
finding a fixed point solution for the RG equation, a sequence of Dirac
operators is determined using (9.58): One starts from a fine lattice, with very
large $\beta$ and a Dirac operator $D$, for the fine lattice where the
couplings are held fixed. The couplings of a new Dirac operator $D'$ for the
coarse lattice are then determined from (9.58). Iterating this step the Dirac
operator at the target $\beta$, which gives rise to the resolution one wants
to work at, may be determined. In each step the coefficients are computed from
(9.58) by minimizing a $\chi^2$ functional which may be chosen as
\begin{equation}
  \chi^2 = \min\sum_i\sum_k
    \left\|\left(D'[U'_{(i)}] - D_{\mathrm{rhs}}[U^{\min}_{(i)}]\right)
      \psi^{\mathrm{rand}}_{(k)}\right\|^2 .
\end{equation}
In this equation $D'[U'_{(i)}]$ is the Dirac operator on the coarse lattice
and $D_{\mathrm{rhs}}[U^{\min}_{(i)}]$ denotes the right-hand side of the RG
equation (9.58). The sum with index $i$ runs over an ensemble of coarse gauge
configurations $U'_{(i)}$, $i = 1, 2, \ldots$, and the corresponding minimal
configurations determined from the FP equation (9.54) are denoted as
$U^{\min}_{(i)}$. The difference of the two Dirac operators then acts on a set
of random vectors $\psi^{\mathrm{rand}}_{(k)}$ labeled by an index $k = 1, 2,
\ldots$. The action of the difference of Dirac operators on the
$\psi^{\mathrm{rand}}_{(k)}$ produces a set of vectors and by $\|\ldots\|$ we
denote the norm of these vectors.

The resulting FP Dirac operator was tested in a series of papers [32--39] and
it was established that the expected properties such as small cutoff effects
can be captured by parameterizations (9.59) and (9.60). Probably even more
important is the fact that the RG equations transport symmetries of the
continuum onto the lattice. For chiral symmetry these properties were
discussed in detail [35, 40, 41]. In particular the FP Dirac operator obeys
the Ginsparg--Wilson equation and the properties which we discussed already in
Chap.~7, such as the correct anomaly, index theorem, follow. We will revisit
the fact that RG transformations transport symmetries onto the lattice in the
next section.

We remark that the parameterization of the Dirac operator (9.60) has been used
also in a different way [38, 42]. Inserting the parameterized Dirac operator
(9.60) into the Ginsparg--Wilson equation (7.29), one can read off a system of
coupled quadratic equations for the coefficients $c^{(k,p)}_{n,m}$. After
truncation this system can be solved numerically and the resulting
coefficients determine the so-called chirally improved Dirac operator. This
operator was tested extensively [39, 43] and was found to have good chiral
properties at a numerical cost considerably lower than necessary for the
overlap operator.

\section{Mapping continuum symmetries onto the lattice}
\label{sec:9.4}

In the last section we have seen that a practical implementation of the RG
approach to full lattice QCD is a rather nontrivial enterprise. We conclude
this chapter by presenting another very important application of the real
space renormalization techniques: We show that the blocking transformation may
be used to map symmetries of the continuum action onto a lattice version of
the symmetry for the corresponding lattice action obtained by blocking. We
discuss this connection between continuum and lattice symmetries here only for
the case of fermions. This case is particularly important, since one of the
symmetries of interest is chiral symmetry -- a long-standing problem of the
lattice formulation. If one maps the (vectorlike) chiral symmetry in the
continuum onto the lattice, one finds that the emerging structure on the
lattice is the Ginsparg--Wilson equation which we already discussed in detail
in Chap.~7. However, the relation between continuum symmetries and their
lattice counterparts is not restricted to chiral symmetries but may be worked
out for general symmetries of a continuum theory.

\subsection{The generating functional and its symmetries}
\label{sec:9.4.1}

Similar to the blocking of free lattice fermions discussed in Sect.~9.2, the
starting point of our analysis is again the equation which defines the Dirac
operator $D$ on the lattice through a blocking transformation from the
fermionic continuum action $S$:
\begin{equation}
  e^{-\bar\psi D\psi} = \int D[\varphi,\bar\varphi]\,
    e^{-(\bar\psi-\bar\varphi^B)B(\psi-\varphi^B) - S_F[\varphi,\bar\varphi]} .
\end{equation}
We denote the lattice fermions by $\psi$, $\bar\psi$ and use vector/matrix
notation for all indices (space--time, color, Dirac). $\varphi$, $\bar\varphi$
are the fermion fields in the continuum which enter the path integral on the
right-hand side. Through integrating over hypercubes which are centered at the
points of the lattice, one constructs from them the blocked fields
$\varphi^B$, $\bar\varphi^B$. These then live on the sites of the lattice and
consequently they have the same indices as the lattice fields $\psi$,
$\bar\psi$, in particular a discrete space--time index.

We stress that (9.62) is understood in a background gauge field which is not
integrated over but also has to be blocked in a suitable way to maintain gauge
covariance (see the last section). The continuum path integral in (9.62) is
defined only formally, but could, e.g., be constructed by a lattice theory in
the continuum limit, as done in the last section. Anyway, here we are only
interested in symmetry properties of the fermion part of the action which is a
bilinear form in the fermion fields. For this analysis a formal definition of
the continuum path integral is sufficient.

Here we consider a more general blocking prescription than the one used in
Sect.~9.2 and allow for a nontrivial blocking kernel $B$ which determines how
the blocked fields $\varphi^B$, $\bar\varphi^B$ and the lattice fields $\psi$,
$\bar\psi$ are mixed. At the moment we leave the blocking kernel $B$
unspecified, but will discuss its role in the end of this section.

The work by Ginsparg and Wilson [44] starts from (9.62) and analyzes its
behavior under a chiral rotation of the lattice fields. Following [45] we use
a slightly different approach and consider the generating functional on the
lattice defined as
\begin{equation}
  W[\bar J,J] = \int D[\psi,\bar\psi]\,e^{\bar\psi J + \bar J\psi - \bar\psi D\psi} ,
\end{equation}
where we have sources $J$ and $\bar J$ coupled to the lattice fermions.
Inserting the exponential of the action from (9.62), we find an expression for
the generating functional through a blocking prescription:
\begin{align}
  W[\bar J,J] &= \int D[\psi,\bar\psi]\,e^{\bar\psi J + \bar J\psi}
    \int D[\varphi,\bar\varphi]\,
    e^{-(\bar\psi-\bar\varphi^B)B(\psi-\varphi^B) - S_F[\varphi,\bar\varphi]} \nonumber\\
  &= \int D[\varphi,\bar\varphi]\,e^{-\bar\varphi^B B\varphi^B - S_F[\varphi,\bar\varphi]}
    \int D[\psi,\bar\psi]\,
    e^{-\bar\psi B\psi + \bar\psi(J+B\varphi^B) + (\bar J+\bar\varphi^B B)\psi} \nonumber\\
  &= \det[B]\,e^{\bar J B^{-1} J}\int D[\varphi,\bar\varphi]\,
    e^{\bar J\varphi^B + \bar\varphi^B J - S_F[\varphi,\bar\varphi]} .
\end{align}
In the last step we have already formally solved the Gaussian integral over
the lattice fields by completing the square in the exponent (compare the
discussion in Sect.~9.2). The resulting expression (9.64) is the lattice
generating functional in terms of a continuum path integral.

We now explore how a symmetry of the continuum fermion action $S_F$ affects
the lattice generating functional $W[\bar J,J]$. In particular we consider a
transformation of the continuum fields:
\begin{equation}
  \varphi \to \varphi' = e^{i\varepsilon T}\varphi , \qquad
  \bar\varphi \to \bar\varphi' = \bar\varphi\,e^{i\varepsilon T} .
\end{equation}
The generators $T$, $T$ of the transformation are matrices acting in Dirac
space. More general transformations are possible as is discussed in [46]. Our
choice implies that the blocked fields $\varphi^B$ and $\bar\varphi^B$
transform in the same way as the original fields, since the blocking from
$\varphi$, $\bar\varphi$ to $\varphi^B$, $\bar\varphi^B$ is a purely scalar
operation, in other words the blocked fields are essentially linear
combinations of the continuum fields. We stress that the generators $T$ and
$\bar T$ are independent transformations, i.e., $\varphi$ and $\bar\varphi$
need not transform in the same way.

We now assume that the transformation (9.65) is a symmetry of the action
\begin{equation}
  S_F[\varphi',\bar\varphi'] = S_F[\varphi,\bar\varphi] .
\end{equation}
One can also evaluate the integral over the continuum fields in the last line
of (9.64) using the transformed variables $\varphi'$, $\bar\varphi'$. Doing so
and exploring the invariance of the action, we obtain
\begin{align}
  &\int D[\varphi',\bar\varphi']\,e^{\bar J\varphi'^B + \bar\varphi'^B J - S_F[\varphi',\bar\varphi']} \nonumber\\
  &= \int D\left[e^{i\varepsilon T}\varphi,\,\bar\varphi\,e^{i\varepsilon T}\right]\,
     \exp\!\left[\bar J e^{i\varepsilon T}\varphi^B
       + \bar\varphi^B e^{i\varepsilon T}J\right. \nonumber\\
  &\qquad \left.- S_F[\varphi,\bar\varphi]\right] \nonumber\\
  &= \left(1 + i\varepsilon A_{TT} + O(\varepsilon^2)\right)
     \int D[\varphi,\bar\varphi]\,
     \exp\!\left[\bar J e^{i\varepsilon T}\varphi^B
       + \bar\varphi^B e^{i\varepsilon T}J\right. \nonumber\\
  &\qquad \left.- S_F[\varphi,\bar\varphi]\right] .
\end{align}
In the last step we have transformed the measure of the continuum path
integral and taken into account that transformation (9.65) could be anomalous
with the anomaly $A_{TT}$ showing up in the Jacobian of the transformation
[47, 48]. Since later we will evaluate all expressions up to $\order{\varepsilon}$,
we have kept only the leading term of the Jacobian. For non-anomalous
transformations $T$, $\bar T$ one has $A_{TT} = 0$.

Inserting result (9.67) back into the expression for the generating functional
(9.64), we find that continuum symmetries (9.65) and (9.66) imply the
following symmetry of the generating functional on the lattice:
\begin{equation}
  W[\bar J,J] = e^{\bar J\left(B^{-1} - e^{i\varepsilon T}B^{-1}e^{i\varepsilon T}\right)J}
    \left(1 + i\varepsilon A_{TT} + O(\varepsilon^2)\right)
    W\left[e^{i\varepsilon T}J,\,\bar J e^{i\varepsilon T}\right] .
\end{equation}
Equation (9.68) summarizes how a continuum symmetry reflects itself in the
lattice generating functional $W[\bar J,J]$ constructed through blocking.

\subsection{Identification of the corresponding lattice symmetries}
\label{sec:9.4.2}

Having analyzed the manifestation of a continuum symmetry in our lattice
generating functional, we now want to identify a corresponding symmetry on the
lattice which also fulfills the symmetry condition (9.68). For that purpose we
consider transformed lattice fields
\begin{equation}
  \psi' = e^{i\varepsilon M}\psi , \qquad
  \bar\psi' = \bar\psi\,e^{i\varepsilon M} .
\end{equation}
The two generators $M$, $M$ for the lattice transformation are not yet known
and we want to identify how they depend on the continuum generators $T$,
$\bar T$ and the blocking kernel $B$. The fact that the transformation should
be a symmetry of the lattice action implies
\begin{equation}
  \bar\psi' D\psi' = \bar\psi\,e^{i\varepsilon M}D\,e^{i\varepsilon M}\psi = \bar\psi D\psi .
\end{equation}
For later use we remark that when expanded in $\varepsilon$ the invariance
condition (9.70) at $\order{\varepsilon}$ implies the commutation relation
\begin{equation}
  M D + D M = 0 .
\end{equation}
As we have done for the continuum expression in the last section, we now
express the generating functional (9.63) in terms of the transformed fields
$\psi'$, $\bar\psi'$ and explore the implications of the symmetry (9.70):
\begin{align}
  W[\bar J,J] &= \int D[\psi',\bar\psi']\,e^{-\bar\psi' D\psi' + \bar\psi' J + \bar J\psi'} \nonumber\\
  &= \det\!\left[e^{i\varepsilon M}\right]\det\!\left[e^{i\varepsilon M}\right]
     \int D[\psi,\bar\psi]\,
     e^{-\bar\psi D\psi + \bar\psi e^{i\varepsilon M}J + \bar J e^{i\varepsilon M}\psi} .
\end{align}
The two Jacobi determinants on the right-hand side come from the
transformation of the measure on the lattice. Using the relation $\det[A] =
\exp(\tr\ln A)$ (see (A.54)) we can expand them as
\begin{equation}
  \det\!\left[e^{i\varepsilon M}\right]\det\!\left[e^{i\varepsilon M}\right]
  = 1 + i\varepsilon\,\tr[M + M] + O(\varepsilon^2) .
\end{equation}
Combining the last two equations we find the symmetry relation for the
generating functional which is implied by lattice symmetry (9.70):
\begin{equation}
  W[\bar J,J] = \left(1 + i\varepsilon\,\tr[M + M] + O(\varepsilon^2)\right)
    W\left[e^{i\varepsilon M}J,\,\bar J e^{i\varepsilon M}\right] .
\end{equation}
Now we compare symmetry condition (9.74) from the lattice transformation to
the corresponding continuum relation (9.68). By setting the two equal we
obtain the equation
\begin{multline}
  e^{\bar J\left(B^{-1} - e^{i\varepsilon T}B^{-1}e^{i\varepsilon T}\right)J}
    \left(1 + i\varepsilon A_{TT} + O(\varepsilon^2)\right)
    W\left[e^{i\varepsilon T}J,\,\bar J e^{i\varepsilon T}\right] \\
  = \left(1 + i\varepsilon\,\tr[M + M] + O(\varepsilon^2)\right)
    W\left[e^{i\varepsilon M}J,\,\bar J e^{i\varepsilon M}\right] ,
\end{multline}
which we can use to identify the lattice transformation $M$, $M$ that matches
the continuum transformation $T$, $\bar T$. The last step is to insert the
explicit form of the generating functional,
\begin{equation}
  W[\bar J,J] = \det[D]\,e^{\bar J D^{-1}J} ,
\end{equation}
which is obtained by directly solving the Gaussian integral (9.63). When
inserting (9.76), (9.75) becomes (we drop the factor $\det[D]$ on both sides)
\begin{multline}
  e^{\bar J\left(B^{-1} - e^{i\varepsilon T}B^{-1}e^{i\varepsilon T}\right)J}
    \left(1 + i\varepsilon A_{TT} + O(\varepsilon^2)\right)
    e^{\bar J e^{i\varepsilon T}D^{-1}e^{i\varepsilon T}J} \\
  = \left(1 + i\varepsilon\,\tr[M + M] + O(\varepsilon^2)\right)
    e^{\bar J e^{i\varepsilon M}D^{-1}e^{i\varepsilon M}J} .
\end{multline}
The last equation holds for arbitrary $\varepsilon$ and arbitrary sources $J$,
$\bar J$. Thus the terms bilinear in $J$ and $\bar J$, as well as the
$\order{\varepsilon}$ terms independent of $J$, $\bar J$, have to match. For
the latter term we conclude
\begin{equation}
  A_{TT} = \tr[M + M] ,
\end{equation}
and thus have identified the lattice counterpart of the continuum anomaly. The
terms bilinear in the sources $J$, $\bar J$ lead to a symmetry relation for the
quark propagator,
\begin{equation}
  B^{-1} - e^{i\varepsilon T}B^{-1}e^{i\varepsilon T} + e^{i\varepsilon T}D^{-1}e^{i\varepsilon T}
  = e^{i\varepsilon M}D^{-1}e^{i\varepsilon M} .
\end{equation}
When expanding in $\varepsilon$ one obtains at $\order{\varepsilon}$
\begin{equation}
  T\left(D^{-1} - B^{-1}\right) + \left(D^{-1} - B^{-1}\right)T
  = M D^{-1} + D^{-1}M .
\end{equation}
This equation is solved by
\begin{equation}
  M = T\left(1 - B^{-1}D\right) , \qquad
  M = \left(1 - DB^{-1}\right)T .
\end{equation}
These are the generators of the lattice symmetry which we wanted to find. They
depend on the continuum generators $T$, $\bar T$, the blocking kernel $B$, and
the lattice Dirac operator $D$. Inserting $M$, $M$ into the symmetry relation
(9.71) one ends up with a nonlinear equation for the lattice Dirac operator
$D$,
\begin{equation}
  T D + D T = D\left[B^{-1}T + T B^{-1}\right]D ,
\end{equation}
which is a generalization of the Ginsparg--Wilson equation [45, 46].

When inserting the explicit form (9.81) into the anomaly equation (9.78) one
obtains the final form for the anomaly
\begin{equation}
  A_{TT} = \tr\left[T + T - T B^{-1}D - D B^{-1}T\right] .
\end{equation}
It is easy to check that when blocking a vector-like continuum theory with a
blocking kernel $B = 2\cdot\mathbb{1}$ and considering a chiral rotation in
the continuum, where $T = \bar T = \gamma_5$, Eq.~(9.82) reduces to the usual
Ginsparg--Wilson relation (7.29). The generators $M$, $M$ of (9.81) are
generators (7.2) of L\"uscher's symmetry and the anomaly assumes the form $A =
-\tr[\gamma_5 D]$, which we have already discussed in Chap.~7.

We have shown that with the help of the blocking transformation we can identify
the mathematical structures that implement chiral symmetry on the lattice, in
particular the Ginsparg--Wilson equation, the chiral transformations on the
lattice, and the form of the anomaly. All of these were developed individually
in Chap.~7, but now follow in one go from a single calculation. Recently the
blocking transformations were explored further with the motivation of
implementing properly Weyl fermions on the lattice as needed for the
electroweak sector of the standard model [46, 49--51]. A central insight [50]
concerns the role of the blocking matrix $B$:
\begin{enumerate}
  \item The blocking matrix $B$ must break all symmetries that are anomalous
        in the target theory.
  \item Other symmetries may be broken by $B$ if it is convenient.
\end{enumerate}
For the case of QCD the choice of $B \propto \mathbb{1}$, which we have used in
the above discussion, is sufficient since only the U(1) axial transformation
is anomalous. For the case of the electroweak sector additional breaking is
required due to the fermion number anomalies. For this case a suitable
blocking matrix was proposed in a particular model [50] and the correct
symmetry pattern was shown to emerge for the corresponding fixed point action.
For the resulting Ginsparg--Wilson equations it is possible to find an explicit
solution [51] using the overlap construction.

\section*{References}

\begin{enumerate}[label={[\arabic*]}]
  \item K.~Symanzik: Nucl.\ Phys.\ B 226, 187 (1983).
  \item K.~Symanzik: Nucl.\ Phys.\ B 226, 205 (1983).
  \item M.~L\"uscher and P.~Weisz: Commun.\ Math.\ Phys.\ 97, 59 (1985).
  \item M.~L\"uscher and P.~Weisz: Erratum, Commun.\ Math.\ Phys.\ 98, 433 (1985).
  \item B.~Sheikholeslami and R.~Wohlert: Nucl.\ Phys.\ B 259, 572 (1985).
  \item M.~L\"uscher, S.~Sint, R.~Sommer, and P.~Weisz: Nucl.\ Phys.\ B 478, 365 (1996).
  \item M.~L\"uscher: in \emph{Probing the Standard Model of Particle
        Interactions}, Proceedings Les Houches 1997, edited by R.~Gupta,
        A.~Morel, E.~de Rafael, and F.~David, Vol.~1 (Elsevier, Amsterdam
        1998).
  \item G.~Curci, P.~Menotti, and G.~Paffuti: Phys.\ Lett.\ B 130, 205 (1983).
  \item G.~Curci, P.~Menotti, and G.~Paffuti: Erratum, Phys.\ Lett.\ B 135, 516 (1984).
  \item R.~Wohlert (unpublished) DESY 87--069 (1987).
  \item M.~L\"uscher and P.~Weisz: Nucl.\ Phys.\ B 479, 429 (1996).
  \item G.~P.~Lepage and P.~B.~Mackenzie: Phys.\ Rev.\ D 48, 2250 (1993).
  \item T.~DeGrand and C.~DeTar: \emph{Lattice Methods for Quantum
        Chromodynamics} (World Scientific, Singapore 2006).
  \item K.~Symanzik: Nucl.\ Phys.\ B 190, 1 (1981).
  \item M.~L\"uscher: Nucl.\ Phys.\ B 254, 252 (1985).
  \item M.~L\"uscher, R.~Narayanan, P.~Weisz, and U.~Wolff: Nucl.\ Phys.\ B 384, 168 (1992).
  \item S.~Sint: Nucl.\ Phys.\ B 421, 135 (1994).
  \item S.~Sint: Nucl.\ Phys.\ B 451, 416 (1995).
  \item M.~L\"uscher et al.: Nucl.\ Phys.\ B 491, 323 (1997).
  \item R.~G.~Edwards, U.~M.~Heller, and T.~R.~Klassen: Nucl.\ Phys.\ (Proc.\ Suppl.) 63, 847 (1998).
  \item K.~Jansen and R.~Sommer: Nucl.\ Phys.\ B 530, 185 (1998).
  \item W.~Bietenholz and U.-J.~Wiese: Nucl.\ Phys.\ B 464, 319 (1996).
  \item K.~G.~Wilson and J.~Kogut: Phys.\ Rep.\ 12C, 76 (1974).
  \item K.~G.~Wilson: Rev.\ Mod.\ Phys.\ 47, 773 (1975).
  \item P.~Hasenfratz: in \emph{Advanced School on Non-Perturbative Quantum
        Field Physics: Peniscola, Spain 2--6 June 1997}, edited by
        M.~Asorey and A.~Dobado, p.~137 (World Scientific, Singapore 1998).
  \item P.~Hasenfratz and F.~Niedermayer: Nucl.\ Phys.\ B 414, 785 (1994).
  \item T.~DeGrand, A.~Hasenfratz, P.~Hasenfratz, and F.~Niedermayer: Nucl.\ Phys.\ B 454, 587 (1995).
  \item T.~DeGrand, A.~Hasenfratz, P.~Hasenfratz, and F.~Niedermayer: Nucl.\ Phys.\ B 454, 615 (1995).
  \item M.~Blatter and F.~Niedermayer: Nucl.\ Phys.\ B 482, 286 (1996).
  \item P.~Hasenfratz et al.: Int.\ J.\ Mod.\ Phys.\ C12, 691 (2001).
  \item C.~Gattringer: Phys.\ Rev.\ D 63, 114501 (2001).
  \item C.~B.~Lang and T.~K.~Pany: Nucl.\ Phys.\ B 513, 645 (1998).
  \item P.~Hasenfratz et al.: Nucl.\ Phys.\ B (Proc.\ Suppl.) 94, 627 (2001).
  \item P.~Hasenfratz et al.: Nucl.\ Phys.\ B (Proc.\ Suppl.) 106, 799 (2002).
  \item P.~Hasenfratz et al.: Nucl.\ Phys.\ B 643, 280 (2002).
  \item S.~Hauswirth: Ph.D.\ thesis, Universit\"at Bern, Switzerland (2002).
  \item T.~J\"org: Ph.D.\ thesis, Universit\"at Bern, Switzerland (2002).
  \item C.~Gattringer: Nucl.\ Phys.\ B (Proc.\ Suppl.) 119, 122 (2002).
  \item C.~Gattringer et al.\ [BGR (Bern-Graz-Regensburg) collaboration]: Nucl.\ Phys.\ B 677, 3 (2004).
  \item P.~Hasenfratz, V.~Laliena, and F.~Niedermayer: Phys.\ Lett.\ B 427, 125 (1998).
  \item P.~Hasenfratz: Nucl.\ Phys.\ B 525, 401 (1998).
  \item C.~Gattringer, I.~Hip, and C.~B.~Lang: Nucl.\ Phys.\ B 597, 451 (2001).
  \item C.~Gattringer et al.: Phys.\ Rev.\ D 79, 054501 (2009).
  \item P.~H.~Ginsparg and K.~G.~Wilson: Phys.\ Rev.\ D 25, 2649 (1982).
  \item C.~Gattringer and M.~Pak: PoS LATTICE2007, 081 (2007).
  \item P.~Hasenfratz, F.~Niedermayer, and R.~von Allmen: JHEP 0610, 010 (2006).
  \item K.~Fujikawa: Phys.\ Rev.\ Lett.\ 42, 1195 (1979).
  \item K.~Fujikawa: Phys.\ Rev.\ D 21, 2848 (1980).
  \item P.~Hasenfratz and M.~Bissegger: Phys.\ Lett.\ B 613, 57 (2005).
  \item P.~Hasenfratz and R.~von Allmen: JHEP 0802, 079 (2008).
  \item C.~Gattringer and M.~Pak: Nucl.\ Phys.\ B 801, 353 (2008).
\end{enumerate}
